% Grade 11 teaching-learning plan: table-based institutional layout.
% Compile with pdfLaTeX. This file is self-contained.
\documentclass[10pt,letterpaper]{article}

\usepackage[landscape,margin=0.5in,bottom=0.62in]{geometry}
\usepackage[T1]{fontenc}
\usepackage[utf8]{inputenc}
\usepackage[scaled=0.94]{helvet}
\renewcommand{\familydefault}{\sfdefault}
\usepackage{microtype}
\usepackage[table]{xcolor}
\usepackage{array,tabularx,multirow}
\usepackage[most]{tcolorbox}
\usepackage{ragged2e}
\usepackage{fancyhdr}

\definecolor{sectionblue}{HTML}{4BACC6}
\definecolor{borderblue}{HTML}{4F81BD}
\definecolor{highlightyellow}{HTML}{FFFF00}

\setlength{\parindent}{0pt}
\setlength{\parskip}{0pt}
\setlength{\tabcolsep}{8pt}
\renewcommand{\arraystretch}{1.12}

\pagestyle{fancy}
\fancyhf{}
\renewcommand{\headrulewidth}{0pt}
\renewcommand{\footrulewidth}{0pt}
\fancyfoot[L]{\colorbox{highlightyellow}{\strut\rmfamily\fontsize{10}{11}\selectfont Versión 08}}

\newcommand{\schoollogo}{%
	\fbox{\parbox[c][0.68in][c]{0.68in}{\centering\scriptsize
			SAN\\BARTOLOMÉ\\[1pt]JHS\\[1pt]LA MERCED}}%
}

\newcommand{\sectionbar}[1]{%
	\begin{tcolorbox}[
		enhanced,
		colback=sectionblue,
		colframe=sectionblue,
		boxrule=0pt,
		sharp corners,
		left=0.32in,
		right=0.08in,
		top=3.5pt,
		bottom=3.5pt,
		before skip=0pt,
		after skip=0pt,
		borderline north={1.15pt}{0pt}{black},
		borderline south={1.15pt}{0pt}{black}
		]
		\color{white}\bfseries\fontsize{12}{14}\selectfont #1
	\end{tcolorbox}%
}

\newtcolorbox{bluebox}[1][]{
	enhanced,
	colback=white,
	colframe=borderblue,
	boxrule=1.15pt,
	sharp corners,
	left=6pt,
	right=6pt,
	top=5pt,
	bottom=5pt,
	before skip=0pt,
	after skip=0pt,
	#1
}

\newtcolorbox{plainbox}[1][]{
	enhanced,
	colback=white,
	colframe=black,
	boxrule=0.45pt,
	sharp corners,
	left=6pt,
	right=6pt,
	top=5pt,
	bottom=5pt,
	before skip=0pt,
	after skip=0pt,
	#1
}

\newcommand{\criterion}[2]{\textbf{C#1.} #2\par}
\newcommand{\field}[2]{\textbf{#1:} #2\par}
\newcommand{\yellowmark}[1]{\colorbox{highlightyellow}{\strut\textcolor{white}{#1}}}

\begin{document}
	\fontsize{9.5}{12}\selectfont
	
	% Institutional header
	\noindent
	\begin{tabularx}{\textwidth}{|>{\centering\arraybackslash}m{0.95in}|>{\centering\arraybackslash}X|>{\centering\arraybackslash}m{1.18in}|}
		\hline
		\begin{minipage}[c][0.86in][c]{0.82in}\centering\schoollogo\end{minipage} &
		{\bfseries\fontsize{12}{14}\selectfont TEACHING-LEARNING PLAN} &
		\begin{minipage}[c][0.86in][c]{1.08in}
			\centering
			\colorbox{highlightyellow}{\bfseries\fontsize{10.5}{12}\selectfont PGF-02-R39}\par
			\vspace{5pt}\hrule\vspace{4pt}
			\textbf{DATE}\par
			\vspace{2pt}
			August 12 --\par November 5, 2026
		\end{minipage}
		\\
		\hline
	\end{tabularx}
	
	\vspace{0.23in}
	
	\sectionbar{1.\quad Identification}
	
	\vspace{0.23in}
	
	\begin{plainbox}
		\field{Area of knowledge}{Global Economics}
		\field{Teacher's name (or Teachers' names)}{Nicolas Lopez}
		\field{Grade}{11}
		\field{Term}{1 (2026-2027)}
	\end{plainbox}
	
	\vspace{0.65in}
	
	\sectionbar{2.\quad Learning Objective}
	
	\vspace{0.22in}
	
	\begin{minipage}[t]{0.28\textwidth}
		\centering\bfseries\fontsize{12}{14}\selectfont LEARNING OBJECTIVE
	\end{minipage}%
	\hfill
	\begin{minipage}[t]{0.70\textwidth}
		\centering\bfseries\fontsize{12}{14}\selectfont ASSESSMENT CRITERIA
	\end{minipage}
	
	\vspace{3pt}
	
	\begin{minipage}[t]{0.28\textwidth}
		\begin{bluebox}[height=2.82in,valign=top]
			\fontsize{9.15}{12.1}\selectfont
			Analyze economic uncertainty through probability models by classifying
			random variables, constructing and normalizing continuous densities,
			calculating and interpreting probabilities for linear, circular,
			exponential and normal models, evaluating model suitability, and
			communicating evidence through cumulative assessment, data-based digital
			modeling and a sourced comparative report.
		\end{bluebox}
	\end{minipage}%
	\hfill
	\begin{minipage}[t]{0.70\textwidth}
		\begin{bluebox}[height=2.82in,valign=top]
			\fontsize{7.95}{10.45}\selectfont
			\criterion{1}{Classify random variables as discrete or continuous; identify support and interpret PMFs, PDFs and CDFs in economic contexts.}
			\criterion{2}{Normalize constant and piecewise-linear densities and calculate interval probabilities geometrically for uniform, triangular, linear and piecewise models.}
			\criterion{3}{Model periodic variables, calculate circular-geometry probabilities, and interpret circular means and concentration.}
			\criterion{4}{Apply the exponential distribution to constant-rate waiting times; calculate probabilities and expected waits, and assess model suitability.}
			\criterion{5}{Use the standard normal distribution and z-table to calculate tail, interval and two-tail probabilities.}
			\criterion{6}{Standardize contextual normal variables and calculate probabilities for practical intervals and thresholds.}
			\criterion{7}{Use inverse-normal reasoning to find percentile cutoffs and central intervals in context.}
			\criterion{8}{Create interactive continuous-density models that display PDFs and shaded below, above and between probabilities.}
			\criterion{9}{Fit and compare probability models to data using parameter estimates, histograms, fitted curves and justified interpretation.}
			\criterion{10}{Produce a sourced comparative report with methods, model probabilities, evidence, conclusions and APA references.}
		\end{bluebox}
	\end{minipage}
	
	\newpage
	
	% Conceptual standards, context, and first learning evidence
	\sectionbar{3. Conceptual standards of the disciplinary knowledge that supports the narrative's development.}
	
	\vspace{0.23in}
	
	\begin{bluebox}
		\begin{center}
			\begin{minipage}{0.89\textwidth}
				\textbf{Core disciplinary knowledge:} random variables and support;
				discrete and continuous models; PMFs, PDFs and CDFs; normalization and total
				probability; Bernoulli, binomial, geometric, Poisson, uniform, triangular,
				linear and piecewise distributions.
			\end{minipage}
		\end{center}
		\textbf{Analytical standards:} geometric probability areas; periodic
		variables, circular means and concentration; exponential waiting times,
		rates, expectation and memorylessness; normal standardization and inverse
		probability; model fit, parameter estimation, interpretation and
		evidence-based communication.
	\end{bluebox}
	
	\vspace{0.24in}
	
	\sectionbar{4.\quad Pre-lesson and Context:
		\yellowmark{\fontsize{10.7}{12}\selectfont Pedagogical actions planned by the teacher to motivate students for the term work proposal}\par
		\hspace*{0.49in}\yellowmark{\fontsize{10.7}{12}\selectfont (presentation of scenarios, narratives, learning objectives and assessment criteria)}}
	
	\vspace{0.22in}
	
	\begin{bluebox}
		Begin with familiar economic and operational variables such as purchases,
		arrivals, waiting times, directions and time of day. Use the compact C1
		classification handout as a diagnostic, then introduce random-variable
		support, probability as area, the term objective, the C1-C10 progression and
		the three learning-evidence blocks. Share materials and instructions through
		OPEN/Teams.
	\end{bluebox}
	
	\vspace{0.23in}
	
	\sectionbar{5.\quad Experience:
		\yellowmark{\fontsize{10.7}{12}\selectfont pedagogical actions that establish the learning objective according to the assessment criteria.}\par
		\hspace*{0.49in}\yellowmark{\fontsize{10.7}{12}\selectfont Links to virtual resources or ways used to share experiences (OPEN, TEAMS, etc.) should be included.}}
	
	\vspace{0.22in}
	
	\renewcommand{\arraystretch}{1.10}
	\noindent
	\begin{tabularx}{\textwidth}{|>{\RaggedRight\arraybackslash}p{2.42in}|>{\RaggedRight\arraybackslash}X|}
		\hline
		\textbf{PEDAGOGICAL ACTION: N1}\par
		\textbf{C1-C4}\par
		\textit{First Learning Evidence}
		&
		\textbf{Learning sequence:} Develop current Practice Activities 1-4 through
		teacher modeling, geometric visualization, guided calculations and individual
		or paired work. Students classify random variables, normalize and interpret
		PDFs, calculate linear and circular areas, and model exponential waiting
		times.\par
		\textbf{Learning evidence:} Cumulative written exam based on the current
		practice activities, assessing C1-C4. Use the compact solved C1 handout and the
		cumulative solution set for feedback, correction and reinforcement.
		\\
		\hline
	\end{tabularx}
	
	\newpage
	
	% Remaining learning evidence and references
	\noindent
	\begin{tabularx}{\textwidth}{|>{\RaggedRight\arraybackslash}p{2.42in}|>{\RaggedRight\arraybackslash}X|}
		\hline
		\textbf{PEDAGOGICAL ACTION: N2}\par
		\textbf{C5-C7 (extension)}\par
		\textit{Second Learning Evidence}
		&
		\textbf{Learning sequence:} Use the retained normal-distribution extension
		resources to connect continuous density work to the standard normal model.
		Students calculate probabilities with a z-table, standardize contextual normal
		variables, and use inverse-normal reasoning to obtain cutoffs and central
		intervals.\par
		\textbf{Learning evidence:} Written exam based on the retained C5-C7 practice
		activities, supported by the standalone normal-table reference and followed by
		targeted correction.
		\\
		\hline
		\textbf{PEDAGOGICAL ACTION: N3}\par
		\textbf{C8-C10 (extension)}\par
		\textit{Third Learning Evidence}
		&
		\textbf{Learning sequence:} Use the retained C8-C10 project brief to build
		interactive continuous-density visualizations, fit candidate probability
		models to uploaded data, display histograms and fitted curves, and compare
		results using documented evidence.\par
		\textbf{Learning evidence:} Digital web project plus an APA-style comparative
		report containing screenshots, model probabilities, parameter estimates,
		interpretation, conclusions, sources and visible records of any AI-assisted
		work.
		\\
		\hline
	\end{tabularx}
	
	\vspace{0.23in}
	
	{\bfseries\fontsize{12}{14}\selectfont References:}
	
	\vspace{3pt}
	
	\begin{bluebox}[height=1.95in,valign=top]
		\textbf{Course map:} README.md - Grade 11 Global Economics, Term 1 probability materials.\par
		\textbf{Practice Activity 1:} \textnormal{\detokenize{G11_T1_C1_practice_activity1.tex}} - random-variable and distribution classification.\par
		\textbf{Practice Activity 2:} \textnormal{\detokenize{G11_T1_C2_practice_activity2.tex}} - geometric probability under continuous PDFs.\par
		\textbf{Practice Activity 3:} \textnormal{\detokenize{G11_T1_C3_practice_activity3.tex}} - periodic and circular probability models.\par
		\textbf{Practice Activity 4:} \textnormal{\detokenize{G11_T1_C4_practice_activity4.tex}} - exponential waiting-time models.\par
		\textbf{Current evidence:} \textnormal{\detokenize{G11_T1_C1C2C3C4_exam0_solution.tex}} - cumulative C1-C4 assessment and solutions.\par
		\textbf{Retained extensions:} \textnormal{\detokenize{old_G11_T3_C5/C6/C7}} practice files and \textnormal{\detokenize{old_G11_T3_L4_C8C9C10_activity.tex}}.
	\end{bluebox}
	
\end{document}